Applying these criteria to all highly cited papers across all cumulative generations, we identified 77 out of 536 papers as being in an obsolescent state. The values of the respective indicators for highly cited papers classified as obsolescent and non-obsolescent are shown in Fig. \ref{fig:arr1}. For papers classified as obsolescent, the mode of the time from publication to citation peak was 4 years, the mode of the time from the peak to substantial citation decay was 7 years, and the mode of the total time from publication to substantial citation decay was 12 years. These findings support the report by P. D. B. Parolo et al., which showed that, in the biological sciences, citation peaks typically occur within five years after publication. In contrast, for papers not classified as obsolescent, the mode of the time from publication to citation peak was markedly different, at 12 years.
Figure \ref{fig:gr7} presents citation trajectories for highly cited papers classified as obsolescent. The time from publication to substantial citation decay ranged from a minimum of 10 years to a maximum of 84 years. The figure also suggests that more recent papers tend to receive higher citation counts immediately after publication.
